\begin{tikzpicture}[font=\scriptsize, >=latex]
% honest limits cartoon: contact + stray cap + current
\node[font=\small\bfseries] at (4.5,4.9) {where breadboards stop being honest};
% contact resistance
\draw[thick, fill=graybg] (0,2.6) rectangle (2.9,4.4);
\node[align=center] at (1.45,3.9) {contact\\resistance};
\node[align=center, font=\tiny, text=slate] at (1.45,3.2) {1--10 $\Omega$ per tie\\(matters at 100 mA+\\or microvolt signals)};
% stray capacitance
\draw[thick, fill=graybg] (3.2,2.6) rectangle (6.1,4.4);
\node[align=center] at (4.65,3.9) {stray\\capacitance};
\node[align=center, font=\tiny, text=slate] at (4.65,3.2) {2--5 pF tie-to-tie\\(kills high-speed,\\fine for LEDs)};
% current limit
\draw[thick, fill=graybg] (6.4,2.6) rectangle (9.3,4.4);
\node[align=center] at (7.85,3.9) {current\\limit};
\node[align=center, font=\tiny, text=slate] at (7.85,3.2) {clips + thin springs\\$\approx$ 1 A max\\(USB still safer)};
\node[align=left, text=slate, text width=9.5cm] at (4.7,1.9)
{The breadboard is not free: every tie point is a small resistor and
a small capacitor glued to your circuit. For everything in this book
--- LEDs, RC fades, sensor filters, logic at a few MHz --- the
extras are below the noise floor. For RF, fast logic, or precision
analog they are the lesson of another book.};
\node[align=left, text=slate, text width=9.5cm] at (4.7,0.5)
{Bridge ahead: when a build earns permanence it graduates to
perfboard (same through-hole parts, soldered) or a PCB (ch18).
Breadboard is for proving, perfboard for keeping.};
\end{tikzpicture}
